%% ─────────────────────────────────────────────────────────────────────────────
\section{Background}
%% ─────────────────────────────────────────────────────────────────────────────

\paragraph{Iterative paraphrasing and authorship.}
In the work done by Tripto et al.~\cite{tripto2024ship} at ACL 2024, the authors ask whether a
text keeps its authorship after repeated paraphrasing. They apply four paraphrasers (ChatGPT,
PaLM2, Pegasus, DIPPER) three times over seven datasets and re-run attribution and
machine-generated-text (MGT) detectors at each generation. Attribution accuracy degrades with
every round, but only aggregate accuracy is reported, so \emph{which} features decay is never
isolated---the gap our work addresses.

\paragraph{Paraphrase as an evasion attack.}
In the work done by Krishna et al.~\cite{krishna2023paraphrasing} at NeurIPS 2023, the concern is
the robustness of AI-text detectors. They train DIPPER, an 11B-parameter discourse-level
paraphraser with controllable lexical and order diversity, and propose a retrieval defense. One
DIPPER pass drives DetectGPT and GPTZero toward chance, but the defense assumes the generating API
logs its outputs.

\paragraph{Zero-shot MGT detection.}
In the work done by Mitchell et al.~\cite{mitchell2023detectgpt} at ICML 2023, the goal is to
detect generated text without labeled data. DetectGPT compares a text's log-likelihood with that of
small perturbations of it, exploiting negative curvature of the model's probability surface. It
beats prior zero-shot baselines but needs white-box likelihoods and collapses under paraphrasing.

\paragraph{Attribution across neural generators.}
In the work done by Uchendu et al.~\cite{uchendu2020authorship} at EMNLP 2020, authorship
attribution is extended from human writers to neural text generators. They benchmark stylometric,
deep-learning and transformer classifiers on three attribution tasks over nineteen generators.
Generators are separable, but only single-pass text is tested, leaving open whether signatures
survive rewriting.

\paragraph{Content-independent author representations.}
In the work done by Rivera-Soto et al.~\cite{riverasoto2021learning} at EMNLP 2021, the question is
whether one authorship representation transfers across domains. They train a Transformer encoder
contrastively on Reddit, Amazon and fan-fiction so that same-author documents embed together. The
embeddings transfer well but remain sensitive to topic overlap, a style--content confound that our
fixed-source design removes.

\paragraph{Measuring semantic retention.}
In the work done by Zhang et al.~\cite{zhang2020bertscore} at ICLR 2020, the problem is evaluating
generated text beyond $n$-gram overlap. BERTScore greedily matches candidate and reference tokens
by contextual-embedding similarity. It tracks human adequacy judgments far better than
BLEU~\cite{papineni2002bleu}, but it is blind to style, so we pair it with explicit stylometry.

\paragraph{Distributional gap between neural and human text.}
In the work done by Pillutla et al.~\cite{pillutla2021mauve} at NeurIPS 2021, the problem is
measuring how far a model's text distribution is from human text. MAUVE compares the two
distributions in an embedding space using divergence frontiers. It correlates well with human
quality judgments, but it scores whole corpora, not individual documents, so it cannot trace one
text's identity across generations.

\paragraph{Watermarking LLM outputs.}
In the work done by Kirchenbauer et al.~\cite{kirchenbauer2023watermark} at ICML 2023, the goal is
to make LLM text verifiably detectable. They bias sampling toward a pseudo-random ``green'' token
list and detect it with a statistical test. The watermark survives light edits but weakens under
paraphrasing, which motivates asking whether the paraphraser leaves a fingerprint of its own (RQ3).

\paragraph{Detection and human judgment.}
In the work done by Ippolito et al.~\cite{ippolito2020automatic} at ACL 2020, the authors study how
decoding strategy affects human and automatic detection of generated text. They benchmark three
sampling methods with expert human raters and fine-tuned BERT classifiers. Sampling that fools
humans makes text easier for classifiers to detect, so both are needed; yet both rely on surface
cues that a paraphrase rewrites.

\paragraph{Visual statistical detection.}
In the work done by Gehrmann et al.~\cite{gehrmann2019gltr} at ACL 2019 (System Demonstrations),
the goal is to help humans spot generated text. GLTR colours each token by its rank under a
language model, exposing unusually predictable text. It raised human detection rates from 54\% to
72\%, but its token-level signal is exactly what a paraphrase replaces.
